\section{Introduction}

Deep learning methods for dynamic economic models fall, broadly, into two families. The first trains a neural network to satisfy the model's equilibrium conditions: Euler equations, first-order conditions, envelope conditions and complementary slackness conditions are written as residuals that vanish at the solution, and the network minimizes their squared sum on simulated states. Deep Equilibrium Nets (DEQNs) \parencite{Azinovic2022DEQN} are the leading example, and they have been applied to integrated assessment models (IAMs) of climate and the economy \parencite{Folini2024ClimateInCE, friedl2023_deepUQ}. The second family, reinforcement learning (RL), maximizes the discounted sum of rewards directly and never writes down an optimality condition \parencite{Sutton2018intro_to_rl, Bertsekas2019RLOC}.

The two families are often presented as close relatives, and they share the structure of a Markov decision process. They differ, however, in what is minimized. A residual-based loss is a consistency criterion: it is zero wherever the hand-derived conditions hold, and it is only as correct as the derivation behind it. If an envelope condition omits a term, the network converges to a wrong fixed point without any sign of failure, since the loss still goes to zero. A reward-based objective, in contrast, is defined by the economic primitives alone (utility, technology and the laws of motion), so there is nothing to derive and nothing to get wrong beyond the model itself.

This paper asks whether a reward-based method can match the accuracy of a correctly specified DEQN on a problem where the exact answer is known. I use the Analytic Climate Economy (ACE) of \textcite{TraegerChristian2023ACE} with the production functions of DICE 2016 and DICE 2023. ACE has a closed-form solution for the consumption rate, the abatement rate, the costates and the social cost of carbon (SCC), which provides an exact benchmark. The motivating example is close at hand: in an earlier version of a DEQN for this model, errors in the envelope condition for capital and double-counted terms in the first-order conditions led the network to a consumption rate between 0.53 and 0.93 instead of the optimal 0.74, and to an abatement rate near zero.

The method is a model-based actor-critic. The critic learns the value of the current policy from the Bellman equation, and the actor maximizes the one-step lookahead, current utility plus the discounted critic value of the next state, by differentiating through the known laws of motion. This is the deterministic policy gradient \parencite{Silver2014DPG} with a known, differentiable model in place of a learned action-value function, as in stochastic value gradients \parencite{Heess2015SVG}. Nothing in the objective is derived by hand.

The main finding is that a plain version of this method fails, and that the reason is economically interesting. With the critic trained by temporal-difference learning alone, the actor learns a reasonable consumption rate but abandons abatement entirely, because the critic attributes part of the growth in value over time to the rising carbon stock and so estimates a \emph{positive} marginal value of atmospheric carbon. The damages from one more gigaton of carbon change the value function by about \(10^{-4}\), far below the precision of a value function fitted to values of order 50. The SCC is a small gradient of a large function, and the policy depends on that gradient, not on the value itself.

Training the critic on its gradient as well as its value resolves the problem. The target gradient is the gradient of the right-hand side of the Bellman equation with respect to the state, computed by automatic differentiation through the known dynamics. This is value-gradient learning \parencite{Werbos1992ADP, Fairbank2012VGL}, a form of Sobolev training \parencite{Czarnecki2017Sobolev}, and it is the reward-based counterpart of the envelope conditions in a DEQN, except that the conditions are generated by the computer rather than derived by the researcher. With each gradient component normalized by a running estimate of its own scale, the actor-critic recovers the analytic consumption rate, abatement path, costates and SCC in both versions of the model. Across three random seeds per model, its SCC is within 1--4\% of the closed-form value and its welfare loss relative to the optimum is below 0.04\% of consumption. A DEQN with correctly derived conditions remains about an order of magnitude more precise, which is the price of learning shadow prices from the curvature of a value function rather than from conditions supplied by the researcher.

The rest of the paper is organized as follows. Section \ref{sec:related} relates the approach to the literature. Section \ref{sec:model} summarizes the model and its analytic solution. Section \ref{sec:method} presents the actor-critic method, and Section \ref{sec:setup} the experimental setup. Section \ref{sec:results} reports the results, Section \ref{sec:discussion} discusses when reward-based and residual-based methods are each appropriate, and Section \ref{sec:conclusion} concludes.
